\documentclass[11pt]{letter}

\usepackage[T1]{fontenc}
\usepackage[utf8]{inputenc}
\usepackage{lmodern}
\usepackage[margin=1in]{geometry}
\usepackage{microtype}
\usepackage{parskip}

\signature{Egberto F. Selerio Jr.\\
Engineering Research and Development for Technology\\
University of San Carlos, Philippines\\
\texttt{egbertoselerio@gmail.com}}

\address{Egberto F. Selerio Jr.\\
Engineering Research and Development for Technology\\
University of San Carlos, Philippines}

\date{\today}

\begin{document}

\begin{letter}{E. N. Pistikopoulos, PhD FIChemE FAIChE FREng\\
Editor-in-Chief\\
\textit{Computers \& Chemical Engineering}\\
Texas A\&M University\\
College Station, Texas, United States}

\opening{Dear Professor Pistikopoulos,}

Please consider our original research article, ``Enforcement of Mass Conservation and Non-Negativity in Statistical Surrogates of Activated Sludge Systems,'' for publication in \textit{Computers \& Chemical Engineering}.

The manuscript addresses a practical weakness in surrogate-assisted process modeling: a fast and accurate prediction may still be unusable if it violates the physical bookkeeping of the process it represents. This problem is especially important for activated sludge modeling, where predicted effluent states are often interpreted directly or passed to downstream screening, optimization, and decision-support workflows. Our study asks whether a single post-prediction feasibility step can enforce two basic requirements--mass conservation and componentwise non-negativity--without redesigning or retraining the underlying statistical surrogate.

To answer this question, we transfer the positivity-preserving conservative projection of Kircher and Votsmeier to steady-state ASM2d-TSN effluent prediction and evaluate it on complete 20-component activated sludge states. The benchmark is intentionally broad: 13 statistical surrogates, eleven in-distribution training-set sizes, nested cross-validation, and two separately simulated out-of-distribution regimes. Every raw surrogate prediction set violated mass conservation, and several accurate raw surrogates also produced negative small-component concentrations. After projection, all $766{,}350$ evaluated predictions conserved mass and remained non-negative, with maximum projected conservation residuals of $2.06\times10^{-13}$ in-distribution and $1.72\times10^{-13}$ out-of-distribution.

The main contribution is therefore not simply another surrogate accuracy comparison. The work shows that predictive accuracy and physical admissibility are distinct requirements. Projection improved most component-level nRMSE comparisons and 20 of 26 out-of-distribution aggregate comparisons, but its aggregate in-distribution effect still depended on the surrogate. This result gives process modelers a clearer way to evaluate statistical surrogates: first determine whether the prediction is accurate enough for the intended use, and separately determine whether the predicted state is physically admissible enough to be interpreted, optimized, or used in subsequent computations.

We believe the article is a strong fit for \textit{Computers \& Chemical Engineering} because it combines process modeling, simulation-based benchmarking, machine learning, and constraint enforcement in a form that is directly relevant to readers working on surrogate modeling, optimization, intelligent systems, and process systems engineering. Beyond the activated sludge case, the study offers a reusable evaluation logic for checking whether data-driven process models produce physically meaningful multicomponent predictions.

The manuscript is original, has not been published previously, and is not under consideration elsewhere. The authors have approved its submission and declare no competing interests. Thank you for your consideration.

\closing{Sincerely,}

\end{letter}
\end{document}
